%%
%% This is file `example_fr.tex',
%% generated with the docstrip utility.
%%
%% The original source files were:
%%
%% coppe.dtx  (with options: `examplefr')
%% 
%% This is one of the demonstration documents of the `coppe' document class:
%% the five main-language demos, the PDF/A build of the full example and the
%% side-by-side cover sheet.  It is generated from coppe.dtx by docstrip; edit
%% the .dtx source, not this file.
%% 
%% Copyright (C) 2008-2026 Vicente Helano F. Batista, George O. Ainsworth Jr.,
%% Geraldo Xexéo and Eduardo Mangeli. Released under the GNU General Public
%% License version 3 (see the COPYING file).
%% 
%% \CheckSum{0}
%% \CharacterTable
%%  {Upper-case    \A\B\C\D\E\F\G\H\I\J\K\L\M\N\O\P\Q\R\S\T\U\V\W\X\Y\Z
%%   Lower-case    \a\b\c\d\e\f\g\h\i\j\k\l\m\n\o\p\q\r\s\t\u\v\w\x\y\z
%%   Digits        \0\1\2\3\4\5\6\7\8\9
%%   Exclamation   \!     Double quote  \"     Hash (number) \#
%%   Dollar        \$     Percent       \%     Ampersand     \&
%%   Acute accent  \'     Left paren    \(     Right paren   \)
%%   Asterisk      \*     Plus          \+     Comma         \,
%%   Minus         \-     Point         \.     Solidus       \/
%%   Colon         \:     Semicolon     \;     Less than     \<
%%   Equals        \=     Greater than  \>     Question mark \?
%%   Commercial at \@     Left bracket  \[     Backslash     \\
%%   Right bracket \]     Circumflex    \^     Underscore    \_
%%   Grave accent  \`     Left brace    \{     Vertical bar  \|
%%   Right brace   \}     Tilde         \~}
%%

\documentclass[french,dsc]{coppe}
\addbibresource{example.bib}
\begin{document}
  \title{Demonstração da classe CoppeTeX 4.1 em português}
  \foreigntitle{CoppeTeX 4.1 class demonstration in English}
  \titlein{french}{D\'emonstration de la classe CoppeTeX 4.1 en fran\c{c}ais}
  \author{Pr\'enom de}{l'Auteur}
  \advisor{Premier}{Directeur}{D.Sc.}{UFRJ}
  \examiner{Premier Examinateur}{D.Sc.}{UFRJ}
  \examiner{Deuxi\`eme Examinateur}{D.Sc.}{UFRJ}
  \department{PESC}
  \date{05}{2026}
  \keyword{Multilingue}
  \keyword{D\'emonstration}
  %% palavras-chave do resumo em idioma estrangeiro (3.1.2.1.5, Anexo F)
  \foreignkeyword{Multilingual}
  \foreignkeyword{Demonstration}
  %% palavras-chave do terceiro resumo, em português (Norma COPPE, secao 3)
  \braziliankeyword{Multilíngue}
  \braziliankeyword{Demonstração}
  \maketitle
  \frontmatter
  \begin{abstract}
    Ceci est le r\'esum\'e principal en fran\c{c}ais. La th\`ese est
    r\'edig\'ee en fran\c{c}ais comme langue principale, avec la
    couverture conserv\'ee en portugais selon la convention
    institutionnelle de l'UFRJ. Lorem ipsum dolor sit amet, consectetur
    adipiscing elit.
  \end{abstract}
  \begin{foreignabstract}
    This is the English foreign abstract. Required in CoppeTeX whichever
    the main language is, for international readability.
  \end{foreignabstract}
  \begin{brazilianabstract}
    Este \'e o resumo em portugu\^es, oferecido pelo ambiente
    \texttt{brazilianabstract} para a leitura da banca brasileira quando
    o idioma principal n\~ao \'e o portugu\^es.
  \end{brazilianabstract}
  % Listes pré-textuelles : seule la table des matières est obligatoire
  % (Manuel UFRJ 3.1.2.1.6). Les listes des figures, des tableaux, des
  % abréviations et des symboles sont facultatives (3.1.2.2) ; celles des
  % quadros, programmes et algorithmes apparaissent quand le travail en
  % contient (norme COPPE, §10). Cette démonstration n'imprime que la table.
  \tableofcontents
  \mainmatter
  \chapter{Introduction}
  Ce chapitre d\'emontre le corps r\'edig\'e en fran\c{c}ais, tandis que
  la couverture et la folha-de-rosto restent en portugais (mod\`ele
  institutionnel UFRJ).
  \section{Le fichier de dépôt : PDF/A}
  Le dépôt à l'UFRJ est exclusivement numérique depuis la résolution CEPG
  n.~246/2023, et le Manuel de l'UFRJ (2.2d) exige le fichier final en
  \textbf{PDF/A}, le format d'archivage à long terme. Le PDF/A est le PDF
  restreint par la norme ISO~19005 pour la conservation : toutes les polices
  sont incorporées, les couleurs déclarent un profil (sRGB) qui les rend
  indépendantes du matériel, les métadonnées sont dans un bloc XMP que lisent
  les dépôts, et rien ne dépend de l'extérieur --- ni contenu externe, ni
  JavaScript, ni chiffrement. Le fichier s'ouvre ainsi de la même façon dans
  plusieurs décennies, sur n'importe quel ordinateur.

  La classe utilise le \textbf{PDF/A-2b} : la partie 2 admet la transparence,
  que la partie 1 interdit et que produisent souvent les boîtes colorées et les
  graphiques, et le niveau \emph{b} garantit une reproduction visuelle fidèle.
  L'option \texttt{pdfa} suffit : la classe charge \texttt{pdfx}, construit les
  métadonnées à partir du titre, de l'auteur et des mots-clés, et incorpore le
  profil de couleur. Activez-la tôt, comme le fait \texttt{example.tex}. Pour
  vérifier, utilisez le validateur veraPDF, profil PDF/A-2b. Le français n'est
  pas une langue de rédaction admise pour une thèse de la COPPE (résolution
  CEPG n.~302/2024, art.~57) : ce document ne fait que démontrer le mécanisme.

  \section{Flux d'écriture}
  pdf\LaTeX{} garde au plus seize fichiers ouverts en écriture, et chaque liste,
  index ou glossaire en occupe un. C'est pourquoi la classe charge
  \texttt{morewrites}. L'option \texttt{semmorewrites} le désactive et gagne
  environ une demi-seconde par passe, mais seulement pour un travail sans index,
  sans index multiples et sans glossaire.
\end{document}
%% 
%%
%% End of file `example_fr.tex'.
